% =============================================================================
% FEHRADVICE EXAMPLE REPORT
% =============================================================================
% Demonstrates usage of fehradvice-report.cls
%
% Compile with: xelatex example-report.tex
% =============================================================================

\documentclass{fehradvice-report}

% =============================================================================
% DOCUMENT METADATA
% =============================================================================

\title{Behavioral Economics Analysis}
\subtitle{Evidence-Based Strategy Report}
\author{FehrAdvice \& Partners AG}
\date{Januar 2026}

\client{Example Client AG}
\project{Behavioral Intervention Design}
\version{1.0}
\classification{Confidential}

% =============================================================================
% DOCUMENT
% =============================================================================

\begin{document}

% Title page
\maketitle

% Table of contents
\tableofcontents
\newpage

% =============================================================================
\section{Executive Summary}
% =============================================================================

Dieser Report demonstriert die Verwendung des FehrAdvice LaTeX-Templates, das auf dem offiziellen Corporate Identity Guide basiert.

\begin{fainfobox}[Key Findings]
\begin{itemize}
    \item Loss Aversion ist der stärkste Treiber ($\lambda = 2.25$)
    \item Default-Effekte erhöhen Compliance um 34\%
    \item Social Norms haben moderierenden Effekt
\end{itemize}
\end{fainfobox}

% =============================================================================
\section{Farbpalette Demonstration}
% =============================================================================

Die FehrAdvice-Farbpalette umfasst vier Primär- und vier Sekundärfarben.

\subsection{Primärfarben}

\begin{center}
\begin{tabular}{llll}
\toprule
\rowcolor{fadarkblue}
\fahead{Farbe} & \fahead{HEX} & \fahead{RGB} & \fahead{Verwendung} \\
\midrule
\rowcolor{fawhite}
Dunkelblau & \texttt{\#024079} & RGB(2,64,121) & Headlines, Akzente \\
\rowcolor{falightgray}
Hellblau & \texttt{\#549EDE} & RGB(84,158,222) & Sekundär \\
\rowcolor{fawhite}
Dunkelgrau & \texttt{\#25212A} & RGB(37,33,42) & Fliesstext \\
\rowcolor{falightgray}
Hellgrau & \texttt{\#F3F5F7} & RGB(243,245,247) & Hintergründe \\
\bottomrule
\end{tabular}
\end{center}

\subsection{Sekundärfarben}

\begin{center}
\begin{tabular}{llll}
\toprule
\rowcolor{fadarkblue}
\fahead{Farbe} & \fahead{HEX} & \fahead{RGB} & \fahead{Verwendung} \\
\midrule
\rowcolor{fawhite}
Flieder & \texttt{\#A1A0C6} & RGB(161,160,198) & Diagramme \\
\rowcolor{falightgray}
Mintgrün & \texttt{\#7EBDAC} & RGB(126,189,172) & Diagramme \\
\rowcolor{fawhite}
Ocker & \texttt{\#DECB3F} & RGB(222,203,63) & Hervorhebungen \\
\rowcolor{falightgray}
Orange & \texttt{\#DE9D3E} & RGB(222,157,62) & Warnungen, CTA \\
\bottomrule
\end{tabular}
\end{center}

% =============================================================================
\section{Typografie}
% =============================================================================

\subsection{Überschriften (Roboto)}

Die Überschriften verwenden \textbf{Roboto} in Dunkelblau:
\begin{itemize}
    \item H1: 28pt Bold
    \item H2: 22pt Bold
    \item H3: 16pt Regular
    \item H4: 13pt Regular
\end{itemize}

\subsection{Fliesstext (Open Sans)}

Der Fliesstext verwendet \textbf{Open Sans Regular} in 11pt mit 16pt Zeilenabstand.
Dies ist ein Beispieltext, der die Lesbarkeit des Fliesstexts demonstriert.
Die Farbe ist Dunkelgrau (\texttt{\#25212A}).

\subsection{Akzente (Playfair Display)}

Für Zitate und besondere Hervorhebungen verwenden wir \faaccent{Playfair Display}:

\begin{quote}
Behavioral Economics zeigt uns, dass Menschen nicht immer rational handeln -- und genau darin liegt die Chance für bessere Entscheidungen.
\end{quote}

% =============================================================================
\section{Boxen und Hinweise}
% =============================================================================

\subsection{Info-Box}

\begin{fainfobox}[Definition: Nudge]
Ein Nudge ist eine sanfte Intervention, die das Verhalten in eine bestimmte Richtung lenkt, ohne Optionen einzuschränken oder wirtschaftliche Anreize zu verändern.
\end{fainfobox}

\subsection{Warnung-Box}

\begin{fawarningbox}[Achtung: Ethische Überlegungen]
Nudges müssen transparent und im besten Interesse der Zielgruppe gestaltet sein. Manipulation ist nicht akzeptabel.
\end{fawarningbox}

\subsection{Erfolg-Box}

\begin{fasuccessbox}[Ergebnis: Intervention erfolgreich]
Die implementierte Default-Option erhöhte die Opt-in-Rate von 23\% auf 78\%.
\end{fasuccessbox}

% =============================================================================
\section{Listen}
% =============================================================================

\subsection{Aufzählung}

Die wichtigsten Behavioral Economics Konzepte:

\begin{itemize}
    \item \textbf{Loss Aversion}: Verluste wiegen schwerer als Gewinne
    \item \textbf{Present Bias}: Präferenz für sofortige Belohnungen
    \item \textbf{Social Norms}: Orientierung am Verhalten anderer
    \item \textbf{Default Effect}: Tendenz zur voreingestellten Option
\end{itemize}

\subsection{Nummerierte Liste}

Implementierungsschritte:

\begin{enumerate}
    \item Zielverhalten definieren
    \item Barrieren analysieren
    \item Intervention designen
    \item Pilotprojekt durchführen
    \item Ergebnisse evaluieren
    \item Skalieren
\end{enumerate}

% =============================================================================
\section{Tabellen}
% =============================================================================

Eine FehrAdvice-konforme Tabelle mit alternierenden Zeilenfarben:

\begin{center}
\begin{tabular}{lrrr}
\toprule
\rowcolor{fadarkblue}
\fahead{Segment} & \fahead{Baseline} & \fahead{Treatment} & \fahead{$\Delta$} \\
\midrule
\rowcolor{fawhite}
Autonomy-Seeking & 23\% & 45\% & +22pp \\
\rowcolor{falightgray}
Guidance-Seeking & 31\% & 78\% & +47pp \\
\rowcolor{fawhite}
Security-Seeking & 28\% & 62\% & +34pp \\
\rowcolor{falightgray}
Connection-Seeking & 19\% & 51\% & +32pp \\
\midrule
\rowcolor{fawhite}
\textbf{Total} & \textbf{25\%} & \textbf{59\%} & \textbf{+34pp} \\
\bottomrule
\end{tabular}
\end{center}

% =============================================================================
\section{Zahlenformatierung}
% =============================================================================

Schweizer Zahlenformat mit Apostroph als Tausendertrennzeichen:

\begin{itemize}
    \item Umsatz: CHF \num{1250000.00}
    \item Teilnehmer:innen: \num{12500}
    \item Effektstärke: $d = \num{0.45}$
\end{itemize}

% =============================================================================
\section{Schlussfolgerungen}
% =============================================================================

\begin{fainfobox}[Empfehlungen]
\begin{enumerate}
    \item \textbf{Sofort}: Default-Optionen implementieren
    \item \textbf{Kurzfristig}: Social Proof Elemente hinzufügen
    \item \textbf{Mittelfristig}: Personalisierte Nudges entwickeln
\end{enumerate}
\end{fainfobox}

% =============================================================================
% APPENDIX
% =============================================================================

\newpage
\appendix
\section{Methodologie}

Details zur verwendeten Methodik finden sich im EBF Framework (Evidence-Based Framework for Economic and Social Behavior).

\section{Referenzen}

\begin{itemize}
    \item Kahneman, D. \& Tversky, A. (1979). Prospect Theory.
    \item Thaler, R. \& Sunstein, C. (2008). Nudge.
    \item Fehr, E. \& Schmidt, K. (1999). Inequity Aversion.
\end{itemize}

\end{document}
